\documentclass[8pt]{beamer}
\usepackage{../config}




\begin{document}
	\subtitle{5to encuentro: Manipulación de datos II}

	\begin{frame}[plain]
		\maketitle
	\end{frame}
	
	
	\begin{frame}[fragile]
		En el encuentro anterior vimos cómo seleccionar columnas (\rscript{select}) y crear nuevas variables (\rscript{mutate}). Hoy completaremos nuestra caja de herramientas de manipulación de datos aprendiendo a trabajar con las filas (filtrarlas y ordenarlas) y, lo más potente de todo: cómo resumir información por grupos.
		
		Para este encuentro, usaremos un dataset simulado de ventas más completo que nos permitirá hacer análisis interesantes.
	\end{frame}
	
	
	
	\begin{frame}[fragile]
\begin{lstlisting}[style=RScript, numbers=none]
library(`:pk:tibble`)

ventas <- tibble(
	vendedor = c("Ana", "Beto", "Carla", "Ana", "Beto", "Carla",
							 "Ana", "Beto", "Carla", "Ana"),
	region = c("Norte", "Sur", "Norte", "Oeste", "Sur", 
						 "Oeste", "Norte", "Sur", "Oeste", "Este"),
	producto = c("Laptop", "Tablet", "Celular", "Tablet", "Laptop", 
							 "Celular", "Laptop", "Tablet", "Celular", "Laptop"),
	monto = c(1200, 450, 800, 400, 1100, 850, 1250, 480, 820, 1180),
	cantidad = c(2, 5, 10, 3, 2, 8, 3, 6, 9, 2),
	satisfaccion = c(4.5, 3.8, 4.2, 4.0, 4.8, 3.5, 4.7, 3.9, 4.1, 4.6)
)

ventas
\end{lstlisting}
\begin{lstlisting}[style=RConsole]
# A tibble: 10 × 6
vendedor region producto monto cantidad satisfaccion
	<chr>    <chr>  <chr>    <dbl>    <dbl>        <dbl>
1 Ana      Norte  Laptop    1200        2          4.5
2 Beto     Sur    Tablet     450        5          3.8
3 Carla    Norte  Celular    800       10          4.2
4 Ana      Oeste  Tablet     400        3          4  
5 Beto     Sur    Laptop    1100        2          4.8
6 Carla    Oeste  Celular    850        8          3.5
7 Ana      Norte  Laptop    1250        3          4.7
8 Beto     Sur    Tablet     480        6          3.9
9 Carla    Oeste  Celular    820        9          4.1
10 Ana      Este   Laptop    1180        2          4.6
\end{lstlisting}
	\end{frame}
	
	
	
	
	
\begin{frame}[fragile]{El verbo \texttt{filter}}
La función \rscript{filter()} nos permite quedarnos solo con las filas (observaciones) que cumplen cierta condición. Es el equivalente a usar \quotation{filtros} en Excel.

Para usar \rscript{filter()}, necesitamos expresar condiciones lógicas que sean VERDADERAS (\rscript{TRUE}) para las filas que queremos mantener.

\titleitem{Operadores de comparaci\'on}
\begin{itemize}
	\item \rscript{>}: Mayor que
	\item \rscript{<}: Menor que
	\item \rscript{>=}: Mayor o igual que
	\item \rscript{<=}: Menor o igual que
	\item \rscript{==}: Igual a 
	\item \rscript{!=}: Distinto de
	\item \lstinline[style=RScript, basicstyle=\normalsize\ttfamily]|%in%|: Pertenece al conjunto (muy útil para no escribir muchos O)	
\end{itemize}
\end{frame}
	
	

\end{document}